\newif\ifgravitystandalone
\let\gravityenddocument\relax
\ifdefined\tnmainflag
\gravitystandalonefalse
\else
\gravitystandalonetrue
\documentclass[11pt]{article}
\usepackage[margin=1in]{geometry}
\usepackage{amsmath,amssymb}
\usepackage{physics}
\usepackage{bm}

\setlength{\emergencystretch}{1.5em}
\raggedbottom

\newcommand{\pubinput}[1]{%
  \IfFileExists{#1}{\input{#1}}{\input{pub/#1}}%
}

\pubinput{physics_constants.tex}

\IfFileExists{tn.aux}{%
  \usepackage{xr}\externaldocument[tn-]{tn}%
}{%
  \IfFileExists{pub/tn.aux}{\usepackage{xr}\externaldocument[tn-]{pub/tn}}{}%
}

\title{Geometric Emergence and the Topological EFE.}
\author{Static Holographic Boundary Theory}
\date{\today}
\def\gravityenddocument{\end{document}}

\begin{document}
\maketitle
\fi

\makeatletter
\providecommand{\manualappendixlabel}[2]{%
  \begingroup
  \def\@currentlabel{#1}%
  \label{#2}%
  \endgroup
}
\newcommand{\gravitymainref}[1]{%
  \@ifundefined{r@#1}{%
    \@ifundefined{r@tn-#1}{\ref{#1}}{\ref{tn-#1}}%
  }{%
    \ref{#1}%
  }%
}
\makeatother

\ifgravitystandalone
\section{Geometric Emergence and the Topological EFE.}\label{app:gravity}\label{sec:gravity-einstein}
\else
\section*{Appendix G (continued): Geometric Emergence and the Topological EFE.}\manualappendixlabel{G}{app:gravity}\manualappendixlabel{G}{sec:gravity-einstein}
\fi

\noindent\textbf{Scope.} This appendix records the gravity-side derivations retained as backup for the anomaly-free $(26,8,312)$ benchmark branch. It is not an independent tuning sector and does not introduce new scan directions: the same modular-$T$ closure data that fix the positive completion residue $c_{\rm dark}$ also fix the positive vacuum loading $\Lambda_{\rm holo}$, the finite horizon register $N$, and the branch-local bulk closure tensor used below as the gravity-side consistency diagnostic. In this appendix the matter-asymmetry bookkeeping is classified as a \emph{Primary Boundary Residue}: the $c_{\rm dark}$ completion sector is read simultaneously as the source of the effective geometric action and as the $B-L$ parity sink that stores the missing conjugate parity bits of the visible projection. The derivation chain below is therefore conditional on the retained SHBT assumptions---static coarse-grained bulk bookkeeping, finite one-copy capacity, and the anomaly-free framing condition $\Delta_{\rm fr}=0$---rather than a claim of a complete dynamical quantum-gravity construction.

\noindent\textbf{Conceptual Scope.} The gravity-side picture used here is static in the deliberately narrow SHBT sense: it describes a ``frozen'' coarse-grained bulk fixed by the completed boundary data, not a time-dependent evaporation model. The finite bit budget $N$ inferred from $\Lambda_{\rm obs}$ is treated throughout as an \emph{Observational Boundary Condition}, not as an internal fit direction. Accordingly, any Page-curve-style identities are used only as bookkeeping constraints on finite information capacity and code-subspace reconstruction. They are not claimed to provide an independent microscopic mechanism for dynamical black-hole evaporation.

\noindent\textbf{Benchmark status.} Unless explicitly stated otherwise, the detuning arguments below are internal stress tests of the retained anomaly-free branch. They show how the displayed closure identities fail under one-step shifts of the benchmark integers or under loss of finite-capacity saturation; they are not presented as no-go theorems for every non-minimal completion outside the one-copy SHBT bookkeeping.

\manualappendixlabel{G.1}{app:curvature-sign-test}
\subsection*{G.1: Bulk Closure Tensor, Holographic Surface Tension, Finite Bit Budget, and the Topological EFE}

On the anomaly-free branch the modular complement required for fixed-parent modular-$T$ closure is the positive completion residue,
\begin{equation}
\label{eq:gravity-closure-residue}
 c_{\rm dark}=\cDarkCompletion\approx \cDarkCompletionFull>0,
\end{equation}
and the same branch data fix the positive vacuum loading
\begin{equation}
\label{eq:gravity-lambda-positive}
 \Lambda_{\rm holo}=+\frac{3\pi}{L_P^2N}=+\lambdaHoloMetersInverseSquared\,\mathrm{m}^{-2}\approx +1.08913883\times10^{-52}\,\mathrm{m}^{-2}>0.
\end{equation}
In the present appendix, the quantity $\Lambda_{\rm holo}$ is not interpreted
as a separately propagating dark-energy fluid. It is the positive
\emph{Holographic Surface Tension} of the completed $(26,8,312)$ branch: the
geometric tension term required to keep the boundary partition function
normalizable on a finite horizon register.
Write the associated de Sitter horizon scale as
\begin{equation}
\label{eq:gravity-horizon-radius}
 R_H
 \equiv
 H_{\Lambda}^{-1}
 =
 \sqrt{\frac{3}{\Lambda_{\rm holo}}},
 \qquad
 A_H=4\pi R_H^2.
\end{equation}
In the frozen-bulk bookkeeping used throughout this appendix, the completed
$SO(10)_{312}$ branch must close on a finite Euclidean $S^3$ of radius $R_H$.
Its holographic screen is the associated horizon $S^2$ of area $A_H$, so the
one-copy capacity is fixed by the saturated holographic count
\begin{equation}
\label{eq:gravity-bit-budget-derived}
 N_{\rm sat}
 \equiv
 \frac{A_H}{4L_P^2}
 =
 \frac{\pi R_H^2}{L_P^2}
 =
 \frac{3\pi}{L_P^2\Lambda_{\rm holo}}.
\end{equation}
This finite-capacity statement is the normalizability condition for the
completed branch partition function. If the horizon size were allowed to run to
infinity, then $N_{\rm sat}\to\infty$, the one-copy character sum would no
longer define a finite-capacity boundary dictionary, and the completed
$SO(10)_{312}$ partition function would cease to be physically admissible as a
normalizable branch state.
These two quantities are not independent outputs of the branch. The same
$c_{\rm dark}$ sector that sources the effective geometric action through the
positive vacuum loading also acts as the storage register for the missing
$B-L$ parity bits of the visible projection. Vacuum loading and parity-sink
bookkeeping are therefore two boundary images of the same completion residue.
This also removes the usual fine-tuning reading of the cosmological term in the
present framework. Here $\Lambda_{\rm holo}$ is not a tiny vacuum counter-term
balanced against a separate large bare contribution; it is the
\emph{Holographic Surface Tension} required to hold the topological moir\'e
pattern of the $(26,8)$ cell on a finite branch-compatible horizon and to bind
the completed $SO(10)_{312}$ partition function to the finite one-copy
register $N_{\rm sat}$. If one were to send $\Lambda_{\rm holo}\to0^+$, then
Eq.~\eqref{eq:gravity-bit-budget-derived} would force $N_{\rm sat}\to\infty$,
so the completed partition function would become divergent rather than
natural. Equivalently, the mass bridge of
Eq.~\eqref{eq:gravity-neutrino-scale-bridge} would collapse the neutrino floor
to zero as $N^{-1/4}\to0$. The positive sign in
Eq.~\eqref{eq:gravity-lambda-positive} is therefore part of the Bianchi lock
itself: finite positive surface tension, finite register, and non-vanishing
neutrino floor are one closure statement viewed in curvature, information, and
mass units.
The same closure is recorded in two branch-fixed matrices that will be used
repeatedly below. Define the \emph{Topological Identity Matrix}
\begin{equation}
\label{eq:gravity-topological-identity-matrix}
 I_{\rm topo}
 \equiv
 \mathbf I_{\rm topo}
 \equiv
 \operatorname{diag}\!\left(
 \,k_{\ell},
 \,k_q,
 \,K
 \right)
 =
 \operatorname{diag}\!\left(
 \,26,
 \,8,
 \,312
 \right).
\end{equation}
This matrix records only the discrete branch labels of the anomaly-free cell.
Separately, define the \emph{Matching Matrix}
\begin{equation}
\label{eq:gravity-topological-identity-basis}
 \mathbf M_{\rm match}
 =
 \operatorname{diag}\!\left(
 \,\kappa_{D_5},
 \,\mathcal R_{\rm GUT},
 \,r_*
 \right)
 =
 \operatorname{diag}\!\left(
 \,\simplexPrefactor,
 \,\frac{8}{28},
 \,\frac{64}{312}
 \right),
 \qquad
 r_*\equiv\frac{\langle\Sigma_{126}\rangle}{\langle\phi_{10}\rangle}=\frac{64}{312}.
\end{equation}
The matching matrix collects the branch-fixed residues that connect the rigid
topological data to the finite-capacity mass bridge, the CKM threshold
closure, and the benchmark scalar alignment. The scalar ratio rescales the
singular spectrum without perturbing the code-subspace isometries. Define the
associated matching defect by
\begin{equation}
\label{eq:gravity-topological-identity-defect}
 \Delta \mathbf M_{\rm match}
 \equiv
 \mathbf M_{\rm match}
 -
 \operatorname{diag}\!\left(
 \,\simplexPrefactor,
 \,\frac{8}{28},
 \,\frac{64}{312}
 \right).
\end{equation}
Then
\begin{equation}
\label{eq:gravity-topological-identity-saturation}
 \Delta \mathbf M_{\rm match}=0
 \qquad\Longleftrightarrow\qquad
 \mathbf M_{\rm match}
 =
 \operatorname{diag}\!\left(
 \,\simplexPrefactor,
 \,\frac{8}{28},
 \,\frac{64}{312}
 \right),
\end{equation}
while $\mathbf I_{\rm topo}$ remains fixed to the discrete branch labels of
Eq.~\eqref{eq:gravity-topological-identity-matrix}. Saturation is therefore
componentwise: $\Delta\mathbf M_{\rm match}=0$ if and only if each matching
residue attains its branch-fixed benchmark value. That is the precise sense in
which the anomaly-free branch is a topological identity equipped with a fixed
matching prescription rather than a continuously deformable family~\cite{georgi1979,babu1993,friedan1987,witten1989}.
Define the branch-closure tensor by
\begin{equation}
\label{eq:gravity-closure-tensor-local}
 \mathcal E_{\mu\nu}
 \equiv
 \frac{c_{\rm dark}M_P^2}{12}\bigl(G_{\mu\nu}+\Lambda_{\rm holo}g_{\mu\nu}\bigr)
 -T^{\rm vis}_{\mu\nu},
 \qquad
 T^{\rm vis}_{\mu\nu}
 \equiv
 T^{(SU(2)_{26})}_{\mu\nu}+T^{(SU(3)_8)}_{\mu\nu}.
\end{equation}
Equation~\eqref{eq:gravity-closure-tensor-local} defines the Bulk Closure Tensor $\mathcal E_{\mu\nu}$. The vanishing of this tensor is the necessary and sufficient condition for a framing-anomaly-free ($\Delta_{\rm fr}=0$) boundary. In the present appendix, $\mathcal E_{\mu\nu}=0$ should be read as a consistency check confirming that the selected RCFT satisfies the standard AdS/CFT isometry requirements needed for code-subspace reconstruction on the anomaly-free branch. The bulk image of modular-$T$ closure is that this tensor can only differ from zero by the same scalar framing defect that measures failure of integer-spectrum closure,
\begin{equation}
\label{eq:gravity-topological-identity-local}
 \mathcal E_{\mu\nu}=\mathcal Q_{\rm iso}\,\Delta_{\rm fr}\,g_{\mu\nu},
\end{equation}
with $\mathcal Q_{\rm iso}\neq 0$ the branch-fixed isotropic anomaly density determined by the parent, visible, and completion central charges. The branch identity therefore reads
\begin{equation}
 \mathcal E_{\mu\nu}=0
 \qquad\Longleftrightarrow\qquad
 \Delta_{\rm fr}=0.
\end{equation}
Bulk diffeomorphism invariance is therefore not an independent postulate in this benchmark appendix: the contracted Bianchi identity is the bulk Ward identity of the same boundary framing condition that sets $\Delta_{\rm fr}=0$.
Equivalently, the effective Einstein system is not an extra dynamical postulate but the covariant restatement of the branch closure,
\begin{equation}
\label{eq:gravity-effective-einstein-local}
 G_{\mu\nu}+\Lambda_{\rm holo}g_{\mu\nu}
 =
 \frac{12}{c_{\rm dark}M_P^2}T^{\rm vis}_{\mu\nu}.
\end{equation}
This is the precise sense in which a four-dimensional bulk description emerges on the anomaly-free branch: the effective Einstein system is the covariant image of the discrete closure data encoded in $\mathbf I_{\rm topo}$, while $\mathbf M_{\rm match}$ enters only through flavor-sector singular-value matching. Equivalently, the branch carries the explicit \textit{Topological EFE Identity}
\begin{equation}
\label{eq:gravity-topological-efe-identity}
 \Delta_{\rm fr}=0
 \qquad\Longleftrightarrow\qquad
 \mathcal E_{\mu\nu}=0
 \qquad\Longleftrightarrow\qquad
 G_{\mu\nu}+\Lambda_{\rm holo}g_{\mu\nu}
 =
 8\pi G_N\,T^{\rm vis}_{\mu\nu},
\end{equation}
with the Newton coupling locked by the completion residue,
\begin{equation}
\label{eq:gravity-newton-lock-local}
 8\pi G_N
 \equiv
 \frac{12}{c_{\rm dark}M_P^2}.
\end{equation}
In that sense Newton's constant is not an independent infrared dial: the gravitational coupling exists on the anomaly-free branch precisely to carry the central-charge deficit $c_{\rm dark}\approx\cDarkCompletion$ left by the visible projection. The parity-storage role of the same residue is isolated separately in Sec.~G.5.

\paragraph{Saturating Information Capacity.}
Evaluating Eq.~\eqref{eq:gravity-bit-budget-derived} on the branch-fixed
vacuum loading of Eq.~\eqref{eq:gravity-lambda-positive} gives the unique
finite register
\begin{equation}
\label{eq:gravity-saturating-information-capacity}
 N_{\rm sat}
 =
 \frac{3\pi}{L_P^2\Lambda_{\rm holo}}
 \approx
 \modularHorizonBits\ \text{bits}.
\end{equation}
This number is not merely an upper bound; it is the unique one-copy storage
capacity on which the completed boundary modular-$T$ spectrum closes. To make
that necessity explicit, write
\begin{equation}
\label{eq:gravity-capacity-detuning-parameter}
 N=N_{\rm sat}+\delta N.
\end{equation}
Holding the completed $(26,8,312)\oplus c_{\rm dark}$ block fixed, a register
detuning changes the surface-tension term by
\begin{equation}
\label{eq:gravity-capacity-lambda-detuning}
 \delta_N\Lambda_{\rm holo}
 =
 \frac{3\pi}{L_P^2}\left(\frac{1}{N_{\rm sat}+\delta N}-\frac{1}{N_{\rm sat}}\right)
 =
 -\frac{3\pi\,\delta N}{L_P^2N_{\rm sat}(N_{\rm sat}+\delta N)}.
\end{equation}
Because the completed block data are held fixed, $\delta_NT^{\rm vis}_{\mu\nu}=0$.
Equation~\eqref{eq:gravity-closure-tensor-local} therefore gives
\begin{equation}
\label{eq:gravity-capacity-closure-detuning}
 \delta_N\mathcal E_{\mu\nu}
 =
 \frac{c_{\rm dark}M_P^2}{12}\,\delta_N\Lambda_{\rm holo}\,g_{\mu\nu}.
\end{equation}
Hence $\delta_N\mathcal E_{\mu\nu}=0$ if and only if $\delta N=0$. If
$\delta N<0$, part of the completed modular-$T$ spectrum lacks storage support
on the one-copy screen; if $\delta N>0$, the excess register is not sourced by
the completed block and the partition sum leaves the one-copy normalizable
sector. Therefore
\begin{equation}
\label{eq:gravity-capacity-iff-closure}
 \mathcal E_{\mu\nu}=0
 \qquad\Longrightarrow\qquad
 N=N_{\rm sat}\approx\modularHorizonBits,
\end{equation}
so $N\approx3.42\times10^{122}$ is the \emph{Saturating Information Capacity}
required for modular-$T$ spectrum closure of the completed block~\cite{friedan1987,witten1989,gko1985}.

\paragraph{Equivalence-Principle proof by discrete detuning.}
Let
\begin{equation}
\label{eq:gravity-discrete-coordinate-vector}
 X\equiv(G_{\rm SM},k_{\ell},k_q,K),
 \qquad
 \delta X_a\equiv\delta_1 X_a\in\{\pm1\}
\end{equation}
 denote the smallest admissible detuning of a single discrete branch coordinate.
For the fixed-parent visible embedding, integer modular-$T$ closure is
equivalent to the descendant-integrality conditions
\begin{equation}
\label{eq:gravity-discrete-branch-integrality}
 \Delta_{\rm fr}=0
 \qquad\Longleftrightarrow\qquad
 \frac{K}{2k_{\ell}}\in\mathbb Z
 \quad\text{and}\quad
 \frac{K}{3k_q}\in\mathbb Z.
\end{equation}
Because these coordinates are discrete rather than continuous, proving
$\delta_1\mathcal E_{\mu\nu}\neq0$ for such one-step moves is the bulk form of
rigidity: there is no infinitesimal deformation direction that preserves the
Einstein identity.
Linearizing Eq.~\eqref{eq:gravity-closure-tensor-local} around the benchmark
cell gives
\begin{equation}
\label{eq:gravity-discrete-closure-variation}
 \delta_1\mathcal E_{\mu\nu}
 =
 \frac{c_{\rm dark}M_P^2}{12}\,\delta_1\Lambda_{\rm holo}\,g_{\mu\nu}
 -\delta_1 T^{\rm vis}_{\mu\nu}.
\end{equation}
For a one-unit leptonic-level detuning with the parent held fixed,
\begin{equation}
\label{eq:gravity-kl-unit-detuning}
 k_{\ell}:26\to26\pm1
 \qquad\Longrightarrow\qquad
 \frac{K}{2k_{\ell}}
 \in
 \left\{\frac{312}{50},\frac{312}{54}\right\}
 =
 \left\{\frac{156}{25},\frac{52}{9}\right\}
 \notin\mathbb Z,
\end{equation}
so the leptonic descendant multiplicity ceases to be integral and therefore
$\delta_1\Delta_{\rm fr}\neq0$. By Eq.~\eqref{eq:gravity-topological-identity-local}
this implies
\begin{equation}
\label{eq:gravity-kl-detuning-nonzero-closure}
 \delta_1\mathcal E_{\mu\nu}
 =
 \mathcal Q_{\rm iso}\,\delta_1\Delta_{\rm fr}\,g_{\mu\nu}
 \neq0.
\end{equation}
For a one-unit quark-level detuning with the parent held fixed,
\begin{equation}
\label{eq:gravity-kq-unit-detuning}
 k_q:8\to8\pm1
 \qquad\Longrightarrow\qquad
 \frac{K}{3k_q}
 \in
 \left\{\frac{312}{21},\frac{312}{27}\right\}
 =
 \left\{\frac{104}{7},\frac{104}{9}\right\}
 \notin\mathbb Z,
\end{equation}
so the descendant multiplicity ceases to be integral and therefore
$\Delta_{\rm fr}\neq0$. By Eq.~\eqref{eq:gravity-topological-identity-local}
this implies
\begin{equation}
\label{eq:gravity-kq-detuning-nonzero-closure}
 \delta_1\mathcal E_{\mu\nu}
 =
 \mathcal Q_{\rm iso}\,\delta_1\Delta_{\rm fr}\,g_{\mu\nu}
 \neq0.
\end{equation}
For a one-unit parent-level detuning with the visible branch held fixed,
\begin{equation}
\label{eq:gravity-parent-unit-detuning}
 K:312\to312\pm1
 \qquad\Longrightarrow\qquad
 K\pm1\not\equiv0\pmod{52}
 \quad\text{and}\quad
 K\pm1\not\equiv0\pmod{24},
\end{equation}
so neither visible descendant multiplicity remains integral and again
$\delta_1\Delta_{\rm fr}\neq0$. Therefore
\begin{equation}
\label{eq:gravity-parent-detuning-nonzero-closure}
 \delta_1\mathcal E_{\mu\nu}
 =
 \mathcal Q_{\rm iso}\,\delta_1\Delta_{\rm fr}\,g_{\mu\nu}
 \neq0.
\end{equation}
For a one-unit visible-current detuning with the branch levels held fixed,
\begin{equation}
\label{eq:gravity-gsm-unit-detuning}
 G_{\rm SM}:15\to15\pm1
 \qquad\Longrightarrow\qquad
 \delta_1\alpha^{-1}_{\rm surf}
 =
 \pm\frac{K}{k_{\ell}+k_q}
 =
 \pm\frac{312}{34}
 =
 \pm\frac{156}{17}
 \neq0,
\end{equation}
so the gauge part of the visible stress tensor shifts,
$\delta_1T^{\rm vis}_{\mu\nu}\supset\delta_1T^{\rm gauge}_{\mu\nu}\neq0$,
and Eq.~\eqref{eq:gravity-discrete-closure-variation} again gives
$\delta_1\mathcal E_{\mu\nu}\neq0$. Thus every one-step move of a
load-bearing coordinate $X_a\in\{G_{\rm SM},k_{\ell},k_q,K\}$ reopens the bulk
closure defect. Because the benchmark branch already satisfies
$\mathcal E_{\mu\nu}^{\rm bench}=0$, the detuned branch obeys
\begin{equation}
\label{eq:gravity-discrete-detuned-closure}
 \mathcal E_{\mu\nu}^{\rm det}
 =
 \delta_1\mathcal E_{\mu\nu}
 \neq 0
 \qquad\text{for every}\qquad
 \delta X_a=\pm1.
\end{equation}
This is the bulk-closure form of the reviewer trap: there is no hidden
compensating direction that can restore $\mathcal E_{\mu\nu}=0$ after a
one-step shift of any load-bearing integer coordinate. Writing the detuned
field equation as
\begin{equation}
\label{eq:gravity-detuned-einstein-identity}
 G_{\mu\nu}+\Lambda_{\rm holo}g_{\mu\nu}
 =
 8\pi G_N\,T^{\rm vis}_{\mu\nu}
 +\mathcal J_{\mu\nu}^{(a)},
 \qquad
 \mathcal J_{\mu\nu}^{(a)}
 \equiv
 \frac{12}{c_{\rm dark}M_P^2}\,\delta_1\mathcal E_{\mu\nu},
\end{equation}
the extra source satisfies $\mathcal J_{\mu\nu}^{(a)}\neq0$ for every
$\delta_1X_a=\pm1$. Because $\mathcal J_{\mu\nu}^{(a)}$ depends on which
coordinate was detuned, the response is no longer universal and cannot be
absorbed into the single branch-fixed Newton coupling of
Eq.~\eqref{eq:gravity-newton-lock-local}. The Einstein identity is therefore
violated and the bulk Equivalence Principle fails under any $\pm1$ discrete
coordinate shift~\cite{friedan1987,witten1989,gko1985}. Equivalently,
\begin{equation}
 \delta X_a=\pm1
 \qquad\Longrightarrow\qquad
 \delta\mathcal E_{\mu\nu}\neq0
 \qquad\Longrightarrow\qquad
 \mathcal J_{\mu\nu}^{(a)}\neq0,
\end{equation}
so freely falling bulk probes no longer couple to a universal geometric source.
Universality of free fall is therefore lost for every one-step detuning of a
load-bearing branch coordinate.
Within this branch bookkeeping, the Bianchi lock and the holographic bit-budget
saturation are the same statement viewed from the bulk and boundary sides. The
vanishing of the Bulk Closure Tensor requires not only $\Delta_{\rm fr}=0$ but
also that the horizon register sit at its saturated value
$N=N_{\rm sat}=\pi R_H^2/L_P^2$; otherwise $\Lambda_{\rm holo}$ is detuned,
the packing residue is no longer fully absorbed by the isotropic vacuum load,
and the bulk Ward identity reopens. In that sense
\begin{equation}
\label{eq:gravity-bianchi-bit-lock}
 \mathcal E_{\mu\nu}=0
 \qquad\Longrightarrow\qquad
 \nabla^{\mu}\mathcal E_{\mu\nu}=0
 \quad\text{with}\quad
 N=N_{\rm sat}\equiv\frac{\pi R_H^2}{L_P^2},
\end{equation}
while Sec.~G.3 spells out the same lock in FRW language.

\paragraph{The Bianchi rebuttal.}
The neighboring fixed-parent cells make the framing/diffeomorphism link explicit. Holding $(k_q,K)=(8,312)$ fixed while detuning the leptonic level away from $k_{\ell}=26$ gives the fractional vacuum modular-$T$ powers
\begin{equation}
\label{eq:gravity-detuned-vacuum-T}
 \frac{K}{2k_{\ell}}-6
 =
 \begin{cases}
 0.24, & k_{\ell}=25,\\[2pt]
 -\dfrac{2}{9}, & k_{\ell}=27,
 \end{cases}
 \qquad\Longrightarrow\qquad
 e^{2\pi i\left(\frac{K}{2k_{\ell}}-6\right)}\neq 1.
\end{equation}
Equivalently, the detuned branches carry non-vanishing framing defects,
\begin{equation}
\label{eq:gravity-detuned-framing-defect}
 \Delta_{\rm fr}(25,8)=0.24,
 \qquad
 \Delta_{\rm fr}(27,8)=\frac{2}{9},
\end{equation}
so the vacuum $T$-power is fractional rather than integer-valued. In bulk language Eq.~\eqref{eq:gravity-topological-identity-local} no longer closes to zero; instead the effective Einstein equations acquire the anomalous isotropic counter-term
\begin{equation}
\label{eq:gravity-detuned-einstein}
 G_{\mu\nu}+\Lambda_{\rm holo}g_{\mu\nu}
 =
 \frac{12}{c_{\rm dark}M_P^2}\Bigl(T^{\rm vis}_{\mu\nu}+T^{\rm anom}_{\mu\nu}\Bigr),
 \qquad
 T^{\rm anom}_{\mu\nu}
 \equiv
 \mathcal Q_{\rm iso}\,\Delta_{\rm fr}\,g_{\mu\nu}.
\end{equation}
The rebuttal point is then sharp: detuning $k_{\ell}$ to $25$ or $27$ does not merely perturb a fit metric; it injects an anomalous counter-term into the covariant bulk equations because the boundary vacuum $T$-power no longer closes integrally. The benchmark branch $(26,8,312)$ is distinguished precisely because $\Delta_{\rm fr}=0$ removes this extra source and leaves the bulk equations counter-term free.

The same branch data also carry the residual packing pressure
\begin{equation}
\label{eq:gravity-delta-topo}
 \delta_{\rm topo}
 \equiv
 1-\kappa_{D_5}
 \approx
 0.01122895,
\end{equation}
which is the non-unit $D_5$ weight-simplex remainder left after visible packing. The same positive remainder appears in the gauge-side closure language of the main text through the benchmark threshold row
\begin{equation}
\label{eq:gravity-alpha-closure}
 \alpha^{-1}(M_{\rm match})
 =
 137.035999\ldots+\delta_{\rm topo}.
\end{equation}
The gravity-side statement is stronger: the same residue is absorbed into the positive dS$_4$ vacuum loading rather than left as a free stress defect.

\paragraph{Gauge-Emergence Cutoff.}
The same branch closure also fixes the ultraviolet gauge-density proxy carried by the holographic screen. With the visible current normalization $G_{\rm SM}=15$, the surface-tension proxy is
\begin{equation}
\label{eq:gravity-surface-gauge-proxy}
 \alpha^{-1}_{\rm surf}
 =
 G_{\rm SM}\frac{K}{k_{\ell}+k_q}
 =
 \frac{15K}{k_{\ell}+k_q}.
\end{equation}
On the retained visible branch $(k_{\ell},k_q)=\benchmarkVisiblePair$ the Diophantine descendant rule forces the parent tower
\begin{equation}
\label{eq:gravity-parent-tower}
 K\in \operatorname{lcm}(2k_{\ell},3k_q)\,\mathbb Z_{>0}
 =\benchmarkParentLevel\,\mathbb Z_{>0},
 \qquad
 K=\benchmarkParentLevel n,
 \quad n\in\mathbb Z_{>0},
\end{equation}
so the surface datum becomes
\begin{equation}
\label{eq:gravity-surface-gauge-tower}
 \alpha^{-1}_{\rm surf}
 =
 15\frac{\benchmarkParentLevel n}{\benchmarkLeptonLevel+\benchmarkQuarkLevel}
=
 \alphaSurfBenchmarkExact\,n.
\end{equation}
The first admissible parent therefore gives
\begin{equation}
\label{eq:gravity-surface-gauge-benchmark}
 K=\benchmarkParentLevel
 \qquad\Longrightarrow\qquad
 \alpha^{-1}_{\rm surf}
 =
 \alphaSurfBenchmarkExact
 \approx
 \alphaSurfBenchmarkDecimal,
\end{equation}
which is the physical benchmark value of the ultraviolet surface gauge datum. By contrast every higher admissible parent lies in the sparse-boundary regime,
\begin{equation}
\label{eq:gravity-sparse-boundary}
 K>312
 \quad\Longrightarrow\quad
 n\ge 2
 \quad\Longrightarrow\quad
 \alpha^{-1}_{\rm surf}
 \ge
 2\alphaSurfBenchmarkExact
 \approx
 275.294118,
\end{equation}
or equivalently
\begin{equation}
\label{eq:gravity-sparse-boundary-direct}
 \alpha_{\rm surf}
 \le
 \frac{17}{4680}
 \approx
 3.63\times10^{-3}
 =
 \frac12\,\alpha_{\rm surf}^{\rm bench}.
\end{equation}
Thus increasing $K$ above the first admissible parent does not create new visible gauge channels; it dilutes the same fifteen visible current slots across a larger parent cycle and makes the boundary gauge network progressively sparser. In the SHBT dictionary bulk gauge fields are the continuum image of these boundary current-current correlators, so the sparse-boundary tower drives the emergent electromagnetic sector toward decoupling. In the formal limit $K\to\infty$ one has $\alpha_{\rm surf}\to0$, and the bulk gauge field vanishes altogether. Already the first higher multiple $K=624$ enters the regime $\alpha^{-1}_{\rm surf}\ge275$, where the boundary gauge density is too weak to support an observationally viable four-dimensional bulk with detectable electromagnetic interactions.

The exclusion of $K>\benchmarkParentLevel$ is therefore geometric rather than merely numerological. A branch with $K>\benchmarkParentLevel$ would produce a bulk in which electromagnetic forces are effectively absent at the emergence scale, so charged matter would fail to organize into the observed four-dimensional sector. The retained parent level is consequently locked not only by framing closure but also by gauge detectability:
\begin{equation}
\label{eq:gravity-gauge-detectability-lock}
 \text{detectable bulk gauge fields}
 \qquad\Longleftrightarrow\qquad
 K=\benchmarkParentLevel
\end{equation}
within the canonical one-copy $SO(10)\supset SU(3)\times SU(2)$ branch used here. In that precise sense the stability of $K$ is anchored by the very existence of measurable gauge fields in the emergent bulk.

Equivalently, define the gauge-support ratio on the admissible parent tower by
\begin{equation}
\label{eq:gravity-gauge-support-ratio}
 \rho_{\rm gauge}(K)
 \equiv
 \frac{\alpha_{\rm surf}(K)}{\alpha_{\rm surf}(\benchmarkParentLevel)}
 =
 \frac{\benchmarkParentLevel}{K}
 =
 \frac{1}{n},
 \qquad
 K=\benchmarkParentLevel n.
\end{equation}
Avoiding the Sparse Boundary onset is the strict support condition
\begin{equation}
\label{eq:gravity-gauge-support-condition}
 \alpha^{-1}_{\rm surf}(K)<2\alphaSurfBenchmarkExact
 \qquad\Longleftrightarrow\qquad
 \rho_{\rm gauge}(K)>\frac12.
\end{equation}
Because the admissible tower has $n\in\mathbb Z_{>0}$, Eq.~\eqref{eq:gravity-gauge-support-condition} implies $n<2$, hence necessarily $n=1$ and therefore
\begin{equation}
\label{eq:gravity-unique-gauge-supporting-parent}
 K=\benchmarkParentLevel=312.
\end{equation}
Thus $K=312$ is the \emph{Unique Gauge-Supporting Parent}: every higher multiple $K\ge624$ lies in the Sparse Boundary regime, where the visible current network is too dilute to sustain an emergent four-dimensional electromagnetic sector~\cite{friedan1987,witten1989,gko1985}.

\manualappendixlabel{G.2}{app:proof-of-mass}
\manualappendixlabel{G.2}{app:holographic-surface-tension}
\manualappendixlabel{G.2}{app:first-law-efe}
\subsection*{G.2: Holographic Surface Tension, Newton Lock, Hubble Friction, and Information Return}

The positive vacuum loading in Eq.~\eqref{eq:gravity-lambda-positive} is the
branch-fixed holographic surface tension of the completed $(26,8,312)$ cell.
Equations~\eqref{eq:gravity-horizon-radius} and
\eqref{eq:gravity-bit-budget-derived} then fix the branch-local de Sitter
register through
\begin{equation}
\label{eq:gravity-horizon-bits}
 N_{\rm holo}\equiv N
 =
 \frac{3\pi}{L_P^2\Lambda_{\rm holo}}
 \approx \maxComplexityCapacity\ \text{bits}.
\end{equation}
The same finite register is the coarse-grained storage capacity for the
$B-L$ parity sink carried by $c_{\rm dark}$. In that sense the de Sitter
horizon count is simultaneously the vacuum-loading register of the geometric
action and the bookkeeping register that stores the missing conjugate parity
bits required for parent-block neutrality.
Combining this with the branch-fixed mass relation
\begin{equation}
\label{eq:gravity-neutrino-scale-bridge}
 m_{\nu}=\kappa_{D_5}M_PN^{-1/4}
\end{equation}
gives
\begin{equation}
\label{eq:gravity-neutrino-scale-bitcount}
 N
 =
 \frac{\kappa_{D_5}^4M_P^4}{m_{\nu}^4}.
\end{equation}
Substituting Eq.~\eqref{eq:gravity-neutrino-scale-bitcount} into
Eq.~\eqref{eq:gravity-horizon-bits}, and using $G_N=L_P^2=M_P^{-2}$ in the
present normalization, yields the \emph{Unity of Scale Identity}
\begin{equation}
\label{eq:gravity-unity-of-scale-identity}
 \Lambda_{\rm holo}
 =
 \frac{3\pi}{\kappa_{D_5}^4}
 G_N m_{\nu}^4.
\end{equation}
In the present appendix Eq.~\eqref{eq:gravity-unity-of-scale-identity} is not
just a derived curiosity. It is the benchmark's \emph{primary internal
falsification metric}: once the anomaly-free branch fixes $\kappa_{D_5}$ and
the observed cosmological loading fixes $\Lambda_{\rm holo}$ through
Eq.~\eqref{eq:gravity-horizon-bits}, the lightest neutrino scale is no longer
free to drift independently. Any empirical failure of
Eq.~\eqref{eq:gravity-unity-of-scale-identity} would be an internal failure of
the branch closure itself rather than a poor phenomenological fit that could be
rescued by retuning hidden parameters. Operationally, once
$\Lambda_{\rm holo}$ and $G_N$ are fixed, Eq.~\eqref{eq:gravity-unity-of-scale-identity}
is the first gravity-side number to check: the lightest neutrino mass must
land in the branch-compatible few-meV regime or the anomaly-free dictionary
fails internally.

\paragraph{Formal proof of the Unity of Scale Identity.}
Starting from Eq.~\eqref{eq:gravity-neutrino-scale-bridge},
\begin{equation}
 m_{\nu}^4
 =
 \kappa_{D_5}^4 M_P^4 N^{-1},
 \qquad\Longrightarrow\qquad
 N
 =
 \frac{\kappa_{D_5}^4 M_P^4}{m_{\nu}^4}.
\end{equation}
Substituting this expression for $N$ into Eq.~\eqref{eq:gravity-horizon-bits}
gives
\begin{equation}
 \Lambda_{\rm holo}
 =
 \frac{3\pi}{L_P^2}
 \frac{m_{\nu}^4}{\kappa_{D_5}^4 M_P^4}.
\end{equation}
Using the branch normalization $G_N=L_P^2=M_P^{-2}$, one has
\begin{equation}
 L_P^2 M_P^4 = G_N^{-1},
\end{equation}
so the previous line becomes
\begin{equation}
 \Lambda_{\rm holo}
 =
 \frac{3\pi}{\kappa_{D_5}^4}
 G_N m_{\nu}^4,
\end{equation}
which is exactly Eq.~\eqref{eq:gravity-unity-of-scale-identity}. Because the
anomaly-free branch has $\kappa_{D_5}>0$, $G_N>0$, and $m_{\nu}>0$, the right-
hand side is strictly positive; the Unity of Scale Identity therefore proves
that the branch-compatible $\Lambda_{\rm holo}$ is a positive surface-tension
term rather than a freely choosable fluid contribution.
The cosmological surface tension, Newton coupling, and neutrino mass floor are
therefore not independent phenomenological dials in this benchmark; they are
the same saturated branch datum expressed in curvature, gravitational, and
defect-mass units.

\paragraph{The scale trap.}
This identity is also the theory's primary internal falsification metric and
simplest internal trap door. If the lightest
neutrino mass were measured significantly above the branch-fixed meV regime---
for example near $10\,\mathrm{meV}$---then
Eq.~\eqref{eq:gravity-neutrino-scale-bitcount} would force a horizon register
$N\propto m_{\nu}^{-4}$ smaller by a factor of roughly
$\sim1.6\times10^2$ than the saturation value fixed by the observed
$\Lambda_{\rm holo}$. Feeding that detuned register back into
Eq.~\eqref{eq:gravity-horizon-bits} would in turn require a surface tension
$\Lambda_{\rm holo}$ larger by the same factor, incompatible with the observed
cosmological constant. By Eq.~\eqref{eq:gravity-capacity-closure-detuning}
this implies $\delta_N\mathcal E_{\mu\nu}\neq0$; by
Eq.~\eqref{eq:gravity-topological-identity-local} the same mismatch is exactly
the reopening of framing closure, i.e. $\Delta_{\rm fr}\neq0$. In that sense a
ten-meV lightest neutrino would not merely shift the benchmark; it would
falsify the anomaly-free branch by forcing a bit budget incompatible with the
observed vacuum loading.

The corresponding geometric surface-tension action is
\begin{equation}
\label{eq:gravity-surface-action}
 S_{\rm surf}[g]
 =
 -\sigma_{\rm holo}\int d^4x\,\sqrt{-g},
 \qquad
 \sigma_{\rm holo}\equiv\frac{\Lambda_{\rm holo}}{8\pi G_N}.
\end{equation}
Varying with respect to the metric gives the effective geometric stress-energy tensor
\begin{equation}
\label{eq:gravity-geom-stress}
 T_{\mu\nu}^{(\mathrm{geom})}
 \equiv
 -\frac{2}{\sqrt{-g}}\frac{\delta S_{\rm surf}}{\delta g^{\mu\nu}}
 =
 -\sigma_{\rm holo}g_{\mu\nu}
 =
 -\frac{\Lambda_{\rm holo}}{8\pi G_N}\,g_{\mu\nu}.
\end{equation}
Equations~\eqref{eq:gravity-effective-einstein-local} and
\eqref{eq:gravity-newton-lock-local} therefore identify the geometric surface
stress with the same completion residue that closes the boundary partition
function. The Newton coefficient is branch-fixed rather than inserted by hand:
\begin{equation}
\label{eq:gravity-newton-lock-restated}
 8\pi G_N
 =
 \frac{12}{c_{\rm dark}M_P^2}
 \approx
 \frac{12}{\cDarkCompletion M_P^2}.
\end{equation}
Gravity is therefore present on the selected branch specifically because the
positive central-charge deficit $c_{\rm dark}$ must be carried by a covariant
bulk channel; without that completion residue there is no Topological EFE to
support the visible stress tensor.

\paragraph{Hubble-Friction Lock.}
The same branch-fixed surface tension also determines the unique isotropic
friction channel that can absorb the packing deficit
$\delta_{\rm topo}=1-\kappa_{D_5}$. On the anomaly-free background,
\begin{equation}
\label{eq:gravity-hubble-lock-preview}
 H_{\Lambda}
 \equiv
 \sqrt{\frac{\Lambda_{\rm holo}}{3}},
 \qquad
 3H_{\Lambda}^2=\Lambda_{\rm holo},
\end{equation}
and the Bianchi lock from Eq.~\eqref{eq:gravity-bianchi-lock-local} with
$\partial_{\nu}\Lambda_{\rm holo}=0$ requires
\begin{equation}
\label{eq:gravity-frw-friction-preview}
 \dot\rho_{\rm vis}+3H_{\Lambda}(\rho_{\rm vis}+p_{\rm vis})=0.
\end{equation}
If one nevertheless postulates a different isotropic damping rate $\Gamma$,
\begin{equation}
\label{eq:gravity-general-friction-preview}
 \dot\rho_{\rm vis}+3\Gamma(\rho_{\rm vis}+p_{\rm vis})=0,
\end{equation}
then subtracting Eq.~\eqref{eq:gravity-frw-friction-preview} gives
\begin{equation}
\label{eq:gravity-friction-uniqueness-preview}
 3(\Gamma-H_{\Lambda})(\rho_{\rm vis}+p_{\rm vis})=0.
\end{equation}
On the matter-carrying branch $\rho_{\rm vis}+p_{\rm vis}\neq0$, and once
$\Delta_{\rm fr}=0$ no second isotropic counter-term survives. Hence the only
solution is
\begin{equation}
\label{eq:gravity-friction-solution-preview}
 \Gamma=H_{\Lambda}.
\end{equation}
Therefore the isotropic Hubble friction $3H_{\Lambda}$ is the unique mechanism
that absorbs the packing deficit $\delta_{\rm topo}$ into the positive vacuum
loading rather than leaving a residual source in $\nabla^{\mu}T^{\rm vis}_{\mu\nu}$.
Section~G.3 gives the expanded contradiction proof in the moir\'e-consistency
language, but the lock itself is already forced by the Topological EFE identity.

Let $B$ be a ball-shaped boundary region, let $H_{\rm mod}(B)$ be its modular Hamiltonian, and let $\xi_B$ be the associated modular-flow vector field that acts as the local boost/isometry at the entangling cut. The first law of entanglement requires
\begin{equation}
\label{eq:gravity-first-law}
 \delta S_{EE}(B)=\delta\langle H_{\rm mod}(B)\rangle.
\end{equation}
For perturbations around the anomaly-free background, the gravitational side of the same identity is the canonical constraint on the bulk Cauchy slice $\Sigma_B$ bounded by the extremal surface of $B$,
\begin{equation}
\label{eq:gravity-first-law-constraint}
 \delta S_{EE}(B)-\delta\langle H_{\rm mod}(B)\rangle
 =
 \frac{6}{c_{\rm dark}M_P^2}
 \int_{\Sigma_B}
 \xi_B^{\nu}\,\delta\mathcal E_{\mu\nu}
 \, d\Sigma^{\mu}.
\end{equation}
Preserving the branch identity means that Eq.~\eqref{eq:gravity-first-law} must hold for every ball $B$ and for every allowed linearized perturbation. The only way for the right-hand side of Eq.~\eqref{eq:gravity-first-law-constraint} to vanish for arbitrary modular-flow generators $\xi_B$ is
\begin{equation}
\label{eq:gravity-linearized-closure}
 \delta\mathcal E_{\mu\nu}=0.
\end{equation}
Substituting the definition of $\mathcal E_{\mu\nu}$ then gives the linearized Einstein field equations on the branch,
\begin{equation}
\label{eq:gravity-linearized-efe}
 \delta G_{\mu\nu}+\Lambda_{\rm holo}\,\delta g_{\mu\nu}
 =
 \frac{12}{c_{\rm dark}M_P^2}\,\delta T^{\rm vis}_{\mu\nu}.
\end{equation}
Thus the entanglement first law is the linearized expression of the same branch closure already encoded in Eq.~\eqref{eq:gravity-effective-einstein-local}. In this appendix, the vanishing of $\delta\mathcal E_{\mu\nu}$ is used as a consistency check: it confirms that the selected RCFT obeys the standard AdS/CFT modular-flow/isometry matching required for the holographic bridge on the anomaly-free branch. Any failure of the linearized Einstein equations would reopen a nonzero $\delta\mathcal E_{\mu\nu}$ and therefore signal a breakdown of that consistency condition.
Equivalently, linearized information return is only compatible with the
branch-fixed background when the finite register remains saturated. A putative
decompactification of the finite $S^3$ background would change
$\Lambda_{\rm holo}$ through Eq.~\eqref{eq:gravity-bit-budget-derived}, violate
the Unity of Scale Identity in Eq.~\eqref{eq:gravity-unity-of-scale-identity},
and reintroduce the same bulk closure defect that the Bianchi lock excludes.

The same information-return statement can be written directly in terms of the bare topological flavor overlap block. On the selected branch the restricted three-channel overlap matrix built from the branch-fixed $SU(2)_{26}$ modular data is
\begin{equation}
\label{eq:gravity-flavor-overlap}
 S_{\rm flav}
 \equiv
 \bigl[S^{(26)}_{q_{\alpha}g_i}\bigr]_{\alpha,i=1}^{3},
 \qquad
 \det S_{\rm flav}\neq0,
 \qquad
 \operatorname{rank}S_{\rm flav}=3.
\end{equation}
Hence the Gram matrix $S_{\rm flav}^{\dagger}S_{\rm flav}$ is positive definite and its polar unitary exists,
\begin{equation}
\label{eq:gravity-flavor-polar}
 U_{\rm flav}
 \equiv
 S_{\rm flav}\bigl(S_{\rm flav}^{\dagger}S_{\rm flav}\bigr)^{-1/2},
 \qquad
 U_{\rm flav}^{\dagger}U_{\rm flav}=\mathbf 1_3.
\end{equation}
If $S_{\rm flav}=W_{\rm flav}\Sigma_{\rm topo}V_{\rm flav}^{\dagger}$ is the
singular-value decomposition of the bare overlap block, then scalar matching is
implemented in the singular basis as
\begin{equation}
\label{eq:gravity-flavor-dressed-overlap}
 S_{\rm dress}
 \equiv
 W_{\rm flav}\bigl(\Sigma_{\rm topo}\mathbf M_{\rm match}\bigr)V_{\rm flav}^{\dagger}.
\end{equation}
Because $\mathbf M_{\rm match}$ is positive diagonal, it rescales only the
singular values and leaves the polar unitary $U_{\rm flav}=W_{\rm flav}V_{\rm flav}^{\dagger}$ unchanged.
If $\det S_{\rm flav}=0$, some nonzero flavor vector would lie in the kernel of $S_{\rm flav}$ and would be projected out before bulk reconstruction. Information loss would then occur already in the flavor sector. Because the benchmark block is rank-$3$, no such null direction exists.

\paragraph{Theorem IV: Holographic Unitary Lock.}
Let
\begin{equation}
\label{eq:gravity-information-return-map}
 \mathcal U_{\rm holo}:\mathcal H_{\rm bulk}^{\rm code}\to\mathcal H_{\partial}^{\rm comp}
\end{equation}
be the branch-fixed bulk-to-boundary reconstruction map from the anomaly-free
code subspace to the completed boundary Hilbert space. On the selected anomaly-free branch, the restriction of $\mathcal U_{\rm holo}$ to the visible three-channel flavor code is the polar unitary $U_{\rm flav}$ extracted from the bare overlap block $S_{\rm flav}$. The Holographic Unitary Lock is therefore secured by the affine level $k_{\ell}=26$ encoded in $S^{(26)}$ and is independent of the Higgs VEV scale, because scalar matching acts only through $\mathbf M_{\rm match}$ on the singular spectrum. Information Return is therefore not an auxiliary postulate but the direct consequence of the branch-fixed unitary lock enforced by the non-singular flavor overlap matrix,
\begin{equation}
\label{eq:gravity-information-return-isometry}
 \langle \mathcal U_{\rm holo}\psi,\mathcal U_{\rm holo}\phi\rangle_{\partial}
 =
 \langle \psi,\phi\rangle_{\rm bulk}
 \qquad\Longleftrightarrow\qquad
 \det S_{\rm flav}\neq0
 \quad
 \text{with}\quad
 \Delta_{\rm fr}=0.
\end{equation}
Equivalently, this theorem should be read as the precise structural requirement for Information
Return in the present framework rather than as a separate dynamical
information-paradox claim: on the anomaly-free branch,
$\mathcal E_{\mu\nu}=0$ together with $\delta\mathcal E_{\mu\nu}=0$ confirms
that the chosen RCFT satisfies the standard AdS/CFT isometry conditions needed
for code-subspace reconstruction.

\paragraph{Proof.}
Because $\det S_{\rm flav}\neq0$, the Hermitian matrix
$S_{\rm flav}^{\dagger}S_{\rm flav}$ is positive definite and admits the
inverse square root appearing in Eq.~\eqref{eq:gravity-flavor-polar}. Its
polar part $U_{\rm flav}$ is therefore unitary, so on the visible code subspace
\begin{equation}
 U_{\rm flav}^{\dagger}U_{\rm flav}=\mathbf 1_3
 \qquad\Longrightarrow\qquad
 \langle U_{\rm flav}\psi,U_{\rm flav}\phi\rangle
 =
 \langle \psi,\phi\rangle.
\end{equation}
Thus no flavor amplitude is projected out and the reconstruction map preserves
inner products. Equation~\eqref{eq:gravity-flavor-dressed-overlap} shows that
any Higgs-sector matching rescales only the positive singular spectrum and does
not alter the polar unitary extracted from the bare $k_{\ell}=26$ overlap
block. This is the direct linear-algebraic content of Theorem~IV: Holographic
Unitary Lock.

If instead $\det S_{\rm flav}=0$, then there exists a nonzero vector
$v\in\mathbb C^3$ with $S_{\rm flav}v=0$. The corresponding reconstructed state
has vanishing boundary norm while the bulk state still has nonzero norm, so the
map loses information already in the flavor sector and cannot be isometric.
The companion background condition is $\Delta_{\rm fr}=0$: by
Eq.~\eqref{eq:gravity-topological-identity-local} this removes the residual
scalar framing source, while Eqs.~\eqref{eq:gravity-first-law-constraint} and
\eqref{eq:gravity-linearized-closure} ensure that no linearized defect reopens
under modular flow. Once the anomaly filter is imposed, the rank-$3$
non-singularity of $S_{\rm flav}$ is precisely what guarantees Information
Return, because the bulk reconstruction map is then isometric on the completed
code subspace.

\manualappendixlabel{G.3}{app:gwb-tilt-corollary}
\manualappendixlabel{G.3}{app:gwb-tilt-requirement}
\manualappendixlabel{G.3}{app:topological-moire-consistency}
\subsection*{G.3: Topological Moir\'e Consistency, Vacuum Loading, and the Bianchi Lock}

The scalar departure from exact scale invariance is the long-wavelength
\emph{Topological Moir\'e Pattern} generated when the discrete $K=312$
boundary lattice is mapped into a continuous four-dimensional horizon
geometry. Writing the scalar mismatch as
\begin{equation}
\label{eq:gravity-moire-scalar}
 \Delta_{\rm Moire}
 \equiv
 1-n_s
 \approx
 0.0352,
 \qquad
 n_s-1\approx -0.0352,
\end{equation}
the observed red tilt is the continuum beat envelope of the branch-fixed
lattice spacing. The same branch also carries the residual packing deficit of
Eq.~\eqref{eq:gravity-delta-topo}, which supplies the tensor-side
\emph{Geometric Friction} of the $D_5$ weight-simplex packing. Defining
\begin{equation}
\label{eq:gravity-geometric-friction}
 \mathcal F_{\rm geom}
 \equiv
 1-\kappa_{D_5}
 =
 \delta_{\rm topo}
 \approx
 0.01122895,
 \qquad
 \kappa_{D_5}=\simplexPrefactor\approx 0.98877,
\end{equation}
the branch-fixed tensor tilt is
\begin{equation}
\label{eq:gravity-tensor-friction}
 n_t\approx -\mathcal F_{\rm geom}=-(1-\kappa_{D_5})\approx -0.0112.
\end{equation}
The positive vacuum loading supplies the exact isotropic damping scale needed
to absorb this packing remainder into the de Sitter background. Writing
\begin{equation}
\label{eq:gravity-hubble-scale}
 H_{\Lambda}\equiv \sqrt{\frac{\Lambda_{\rm holo}}{3}},
 \qquad
 3H_{\Lambda}^2=\Lambda_{\rm holo},
\end{equation}
the Bianchi identity for the closure tensor is
\begin{equation}
\label{eq:gravity-bianchi-lock-local}
 \nabla^{\mu}\mathcal E_{\mu\nu}
 =
 \frac{c_{\rm dark}M_P^2}{12}\,\partial_{\nu}\Lambda_{\rm holo}
 -\nabla^{\mu}T^{\rm vis}_{\mu\nu}
 =0.
\end{equation}
Because $\partial_{\nu}\Lambda_{\rm holo}=0$ on the completed $(26,8,312)$
branch, Eq.~\eqref{eq:gravity-bianchi-lock-local} reduces to the visible
stress-tensor conservation law
\begin{equation}
\label{eq:gravity-visible-stress-conservation}
 \nabla^{\mu}T^{\rm vis}_{\mu\nu}=0.
\end{equation}
Equivalently, $\nabla^{\mu}T_{\mu\nu}=0$ is a structural requirement of the
$(26,8,312)$ branch, with $T_{\mu\nu}$ understood as the branch-completed
stress tensor whose visible projection is $T^{\rm vis}_{\mu\nu}$.
On the isotropic anomaly-free branch $\partial_{\nu}\Lambda_{\rm holo}=0$, so
the only allowed residue is the scalar friction $\mathcal F_{\rm geom}$ of
Eq.~\eqref{eq:gravity-geometric-friction}; no open vector stress survives.
Equivalently, in the effective FRW rewriting the continuity equation takes the
form
\begin{equation}
\label{eq:gravity-frw-friction}
 \dot\rho_{\rm vis}+3H_{\Lambda}(\rho_{\rm vis}+p_{\rm vis})=0,
\end{equation}
with $H_{\Lambda}$ fixed entirely by the positive vacuum loading
$\Lambda_{\rm holo}$. Because the anomaly-free branch supplies no second
isotropic damping scale, Eq.~\eqref{eq:gravity-frw-friction} can be read as an
exact cancellation statement for the lifted packing remainder. Writing a
putative mismatch as
\begin{equation}
\label{eq:gravity-frw-lock-mismatch}
 \dot\rho_{\rm vis}+3H_{\Lambda}(\rho_{\rm vis}+p_{\rm vis})
 =
 \Delta_{\rm lock}(\rho_{\rm vis}+p_{\rm vis}),
\end{equation}
the anomaly-free branch requires $\Delta_{\rm lock}=0$. Here
$\Delta_{\rm lock}$ is the uncancelled FRW lift of the branch-fixed packing
deficit $\delta_{\rm topo}=1-\kappa_{D_5}$. If the Hubble friction $3H_{\Lambda}$
did not exactly match the geometric cancellation required by
$\delta_{\rm topo}$, then $\Delta_{\rm lock}\neq0$ and therefore
\paragraph{Geometric Friction Identity.}
To make the uniqueness claim explicit, replace the de Sitter damping coefficient
by a general isotropic rate $\Gamma$ and write
\begin{equation}
\label{eq:gravity-general-friction}
 \dot\rho_{\rm vis}+3\Gamma(\rho_{\rm vis}+p_{\rm vis})=0.
\end{equation}
Subtracting Eq.~\eqref{eq:gravity-frw-friction} gives the identity
\begin{equation}
\label{eq:gravity-friction-identity}
 3(\Gamma-H_{\Lambda})(\rho_{\rm vis}+p_{\rm vis})
 =
 \Delta_{\rm lock}(\rho_{\rm vis}+p_{\rm vis}).
\end{equation}
On the exact level-matched branch,
\begin{equation}
\label{eq:gravity-level-matching-local}
 K=2I_{\ell}^{\star}k_{\ell}=3I_q^{\star}k_q
 \qquad\Longrightarrow\qquad
 (k_{\ell},k_q,K)=(26,8,312),
\end{equation}
so $\Delta_{\rm fr}=0$ and the only surviving scalar residue is the packing
deficit $\delta_{\rm topo}=\mathcal F_{\rm geom}$. No second isotropic
counter-term is allowed once the framing anomaly is removed. Therefore the
Bianchi lock requires $\Delta_{\rm lock}=0$, and on the matter-carrying visible
branch $\rho_{\rm vis}+p_{\rm vis}$ is not identically zero. Equation
\eqref{eq:gravity-friction-identity} then has the unique solution
\begin{equation}
\label{eq:gravity-unique-friction}
 \Gamma=H_{\Lambda}.
\end{equation}
Hence the Hubble friction $3H_{\Lambda}$ is the unique isotropic sink that
absorbs $\delta_{\rm topo}$ into the positive vacuum loading. Detuning
$k_{\ell}$ away from $26$ reopens $\Delta_{\rm fr}\neq0$, activates the
anomalous counter-term in Eq.~\eqref{eq:gravity-detuned-einstein}, and
therefore destroys this uniqueness.
Thus,
\begin{equation}
\label{eq:gravity-visible-stress-mismatch}
 \nabla^{\mu}T^{\rm vis}_{\mu 0}
 =
 \dot\rho_{\rm vis}+3H_{\Lambda}(\rho_{\rm vis}+p_{\rm vis})
 =
 \Delta_{\rm lock}(\rho_{\rm vis}+p_{\rm vis})
 \neq 0,
\end{equation}
so the visible stress tensor would fail to be conserved. The term $3H_{\Lambda}$
is therefore the exact Hubble friction that absorbs the packing deficit
$\delta_{\rm topo}\approx0.0112$ into the positive vacuum loading rather than
leaving a Bianchi-violating source term. The same residual therefore reappears as the finite tensor drag $n_t\approx-\delta_{\rm topo}$,
while the scalar sector records the longer-wavelength moir\'e beat
$\Delta_{\rm Moire}=1-n_s$. Together they express the same branch-fixed bulk
consistency: vacuum loading, geometric friction, and the Bianchi lock are one
and the same closure statement viewed at curvature, transport, and tilt level.
At the level of the disclosed branch residues this statement is an exact
Topological Identity: $\mathcal F_{\rm geom}\approx0.01122895$ is the scalar remainder
that the continuity equation must absorb, and the tensor tilt
$n_t\approx-\mathcal F_{\rm geom}\approx-0.0112$ is the tensor-side imprint of
that same conservation law. If either relation failed, the packing deficit
would reappear as $\nabla^{\mu}T_{\mu\nu}\neq0$ and the Bulk Closure Tensor
$\mathcal E_{\mu\nu}$ would no longer vanish on the selected branch.

\paragraph{Gauge-Emergence Exclusion Principle.}
The Bianchi lock also carries a gauge-side exclusion statement. The
Gauge-Emergence Cutoff
\begin{equation}
\label{eq:gravity-gauge-emergence-cutoff-local}
 K_{\rm GE}\equiv\benchmarkParentLevel=312
\end{equation}
is a physical exclusion principle rather than a bookkeeping convention. Because
the admissible parent tower is $K=\benchmarkParentLevel n$, any
\begin{equation}
\label{eq:gravity-sparse-boundary-cutoff-local}
 K\ge624
 \qquad\Longrightarrow\qquad
 n\ge2
 \qquad\Longrightarrow\qquad
 \alpha^{-1}_{\rm surf}\ge 2\alphaSurfBenchmarkExact,
 \quad
 \alpha_{\rm surf}\le\frac{17}{4680},
\end{equation}
so the branch enters the \emph{Sparse Boundary} regime already identified in
Sec.~G.1. In that regime the visible current-current network is too dilute to
sustain an emergent four-dimensional electromagnetic sector: the gauge density
falls below the branch-local support threshold, and the bulk no longer carries
observationally viable electromagnetism. Thus $K=312$ is not merely the first
allowed parent level; it is the unique parent level that is simultaneously
Bianchi-locked and gauge-supporting, while every $K\ge624$ is physically
excluded by sparse-boundary gauge dilution~\cite{friedan1987,witten1989,gko1985}.

\paragraph{Forbidden-Void Proposition.}
A vanishing cosmological constant ($\Lambda = 0$) is topologically forbidden
on the $(26,8,312)$ branch. It necessitates a divergent bit-budget
($N\to\infty$), rendering the boundary partition function physically
non-normalizable. Through Eq.~\eqref{eq:gravity-neutrino-scale-bridge} the same
limit also collapses the neutrino floor~\cite{friedan1987,witten1989,gko1985}. Eq.~\eqref{eq:gravity-bit-budget-derived} makes this implication
explicit on the one-copy lattice. In implication form,
\begin{equation}
\label{eq:gravity-forbidden-void-chain}
 \Lambda_{\rm holo}=0
 \Longrightarrow
 H_{\Lambda}=0
 \Longrightarrow
 N_{\rm sat}=\infty
 \Longrightarrow
 Z_{\partial}^{(26,8,312)}\ \text{non-normalizable}
 \Longrightarrow
 m_{\nu}=\kappa_{D_5}M_PN^{-1/4}=0.
\end{equation}
This proposition is the boundary-normalizability form of the Bianchi lock:
$\Lambda_{\rm holo}=0$ is not a soft limit inside the admissible branch family,
but an exit from the finite-capacity Hilbert sector of the completed
$(26,8,312)$ block.
The proof-by-contradiction is then immediate. Assume that the completed
boundary modular-$T$ matrix still closed on the anomaly-free branch while
$\Lambda_{\rm holo}=0$. Then Eq.~\eqref{eq:gravity-hubble-scale} gives
$H_{\Lambda}=0$, and Eq.~\eqref{eq:gravity-bit-budget-derived} gives
$N_{\rm sat}\to\infty$, so the finite one-copy register decompactifies and the
partition function leaves the normalizable finite-capacity sector. The mass
bridge in Eq.~\eqref{eq:gravity-neutrino-scale-bridge} then forces
$m_{\nu}=\kappa_{D_5}M_PN^{-1/4}\to0$, i.e. the neutrino floor collapses. The
FRW friction channel disappears, but the branch-fixed packing residue
$\delta_{\rm topo}=1-\kappa_{D_5}>0$ remains. Equation
\eqref{eq:gravity-frw-lock-mismatch} would then force
$\Delta_{\rm lock}\neq0$, hence
\begin{equation}
\label{eq:gravity-bianchi-lock-contradiction}
 \nabla^{\mu}T^{\rm vis}_{\mu\nu}\neq0
 \qquad\Longrightarrow\qquad
 \nabla^{\mu}\mathcal E_{\mu\nu}\neq0,
\end{equation}
in contradiction with the Bianchi lock in
Eq.~\eqref{eq:gravity-bianchi-lock-local}. Therefore $\Lambda_{\rm holo}=0$
cannot satisfy both conditions simultaneously: it would require an infinite
bit-budget and thereby violate the finite-capacity premise of the
$(26,8,312)$ lattice. The only branch-compatible solution for which the
completed boundary modular-$T$ matrix closes and the bulk Bianchi identity is
satisfied is the positive vacuum loading already fixed in
Eq.~\eqref{eq:gravity-lambda-positive},
\begin{equation}
 \Lambda_{\rm holo}>0.
\end{equation}

\paragraph{Trap note.}
This is the Bianchi-lock form of the same scale trap derived in
Sec.~G.2. If the lightest neutrino mass were measured significantly
outside the predicted few-meV regime, the resulting bit-budget $N$
would be incompatible with the observed expansion rate, forcing a
non-vanishing framing anomaly $\Delta_{\rm fr}\neq0$ and collapsing the
holographic dictionary~\cite{friedan1987,witten1989,gko1985}.

\manualappendixlabel{G.4}{app:holographic-computational-gradients}
\manualappendixlabel{G.4}{app:clock-skew-derivation}
\subsection*{G.4: Holographic Computational Gradients}

The anomaly-free branch carries a finite horizon register,
$N=3\pi/(L_P^2\Lambda_{\rm holo})$, and that finite register is not only a
storage bound but also a finite update budget. Let $N_{\rm load}(t)$ denote the
cumulative number of boundary bits that have been loaded onto the completion
sector by cosmic time $t$, and define the normalized completion fraction
\begin{equation}
\label{eq:gravity-loading-fraction}
 f_{\rm load}(t)
 \equiv
 \frac{N_{\rm load}(t)}{N},
 \qquad
 0\leq f_{\rm load}(t)\leq 1.
\end{equation}
The physical interpretation of the cosmological expansion is then explicit:
\emph{The expansion of the universe is the loading speed of new boundary bits
onto the $c_{\rm dark}$ stabilizer.}
In the FRW rewriting this means that the late-time expansion rate is the
coarse-grained write speed of the completion fraction,
\begin{equation}
\label{eq:gravity-loading-hubble}
 H_{\rm eff}(t)
 =
 H_{\Lambda}
 +
 \frac{1}{3}\frac{d f_{\rm load}}{dt}.
\end{equation}

The clock-skew statement follows because the loading process is read on the
past light cone, where shells at different redshift do not share the same
proper write clock. If $\tau_{\rm lock}$ is the proper time of a local
boundary-update cycle, then the observer-frame and source-frame ticks are
related by the usual redshift factor,
\begin{equation}
\label{eq:gravity-clock-skew}
 d\tau_{\rm lock} = \frac{dt}{1+z}.
\end{equation}
Therefore the loading law becomes
\begin{equation}
\label{eq:gravity-redshift-dependent-hubble}
 H_0(z)
 =
 H_{\Lambda}
 +
 \frac{1}{3(1+z)}\frac{d f_{\rm load}}{d\tau_{\rm lock}}.
\end{equation}
Equation~\eqref{eq:gravity-redshift-dependent-hubble} is the gravity-side
clock-skew derivation of a redshift-dependent Hubble intercept: once the
boundary register is finite, the same loading process is sampled at different
stages of the write cycle by low-$z$ and high-$z$ observations. A strictly
redshift-independent $H_0$ would require either
$df_{\rm load}/d\tau_{\rm lock}=0$ (no continuing boundary loading) or the
decompactification limit $N\to\infty$ in which the discrete update cycle
ceases to be physical. Neither possibility is compatible with the anomaly-free
branch, which carries a finite register and a nonzero positive vacuum loading.
The local and CMB Hubble determinations therefore probe different segments of
the same finite holographic computation rather than different fluids.

\manualappendixlabel{G.5}{app:gravity-parity-storage}
\subsection*{G.5: Gravitational Storage of B-L Parity}

The visible branch does not erase its missing antimatter complement; it stores
that complement as a topological conjugate in the same geometric sector that is
measured macroscopically by the completion residue $c_{\rm dark}$. In branch
language the visible baryon excess is admissible only because the completed
parent block remains neutral,
\begin{equation}
\label{eq:gravity-parent-neutrality}
 (B-L)_{\rm vis}+(B-L)_{\rm dark}=0,
 \qquad
 (B-L)_{\rm dark}=-(B-L)_{\rm vis}.
\end{equation}
Here $(B-L)_{\rm dark}$ is not a second thermal fluid and not a separately
propagating antimatter bath. It is the non-propagating parity sink of the
geometric completion sector: the missing antimatter of the visible projection
is stored as a topological conjugate in the $c_{\rm dark}$ branch residue.

This is why the same Holographic Stabilizer can play two roles without
contradiction. The residue $c_{\rm dark}$ sources the positive vacuum loading
$\Lambda_{\rm holo}$ in the Topological EFE, and the same residue stores the
conjugate $B-L$ bits needed to keep the projected boundary neutral. The
gravitational field and the parity sink are therefore two macroscopic images of
one and the same completion sector.

The finite de Sitter register of Eq.~\eqref{eq:gravity-horizon-bits} gives the
storage capacity for this parity bookkeeping. In that sense gravitational
storage is literal within the benchmark logic: the branch does not hide the
missing antimatter in an untracked particle sector, but in the geometric
closure data of the completed boundary block. The bulk is neutral because the
same $c_{\rm dark}$ sector that closes $\mathcal E_{\mu\nu}$ also stores the
conjugate parity information required by Eq.~\eqref{eq:gravity-parent-neutrality}.

\ifgravitystandalone
\begin{thebibliography}{99}
\bibitem{georgi1979}
H.~Georgi and C.~Jarlskog,
``A new lepton-quark mass relation in a unified theory,''
\emph{Phys. Lett. B}~86 (1979) 297--300.

\bibitem{babu1993}
K.~S.~Babu and R.~N.~Mohapatra,
``Predictive neutrino spectrum in minimal $SO(10)$ grand unification,''
\emph{Phys. Rev. Lett.}~70 (1993) 2845--2848.

\bibitem{friedan1987}
D.~Friedan and S.~Shenker,
``The Analytic Geometry of Two-Dimensional Conformal Field Theory,''
\emph{Nucl. Phys. B}~281 (1987) 509--545.

\bibitem{witten1989}
E.~Witten,
``Quantum field theory and the Jones polynomial,''
\emph{Commun. Math. Phys.}~121 (1989) 351--399.

\bibitem{gko1985}
P.~Goddard, A.~Kent, and D.~Olive,
``Virasoro algebras and coset space models,''
\emph{Phys. Lett. B}~152 (1985) 88--92.
\end{thebibliography}
\fi

\gravityenddocument
